
\documentclass{article}
\usepackage{colortbl}
\usepackage{makecell}
\usepackage{multirow}
\usepackage{supertabular}

\begin{document}

\newcounter{utterance}

\centering \large Interaction Transcript for game `eqbench', experiment `eq\_bench', episode 57 with reflexion{-}llama{-}3.1{-}8b{-}instruct.
\vspace{24pt}

{ \footnotesize  \setcounter{utterance}{1}
\setlength{\tabcolsep}{0pt}
\begin{supertabular}{c@{$\;$}|p{.15\linewidth}@{}p{.15\linewidth}p{.15\linewidth}p{.15\linewidth}p{.15\linewidth}p{.15\linewidth}}
    \# & \multicolumn{2}{c}{Player} && \multicolumn{2}{c}{Game Master} \\
    \hline

    \theutterance \stepcounter{utterance}  
    & & & \multicolumn{4}{p{0.6\linewidth}}{
        \cellcolor[rgb]{0.9,0.9,0.9}{
            \makecell[{{p{\linewidth}}}]{
                \texttt{\tiny{[P1$\langle$GM]}}
                \texttt{You have attempted this same task before and did not succeed. These are the plans you wrote after those attempts.} \\
\texttt{Trial \#1: New plan: I overestimated Claudia's emotional state and didn't consider the context of the conversation. Upon re{-}reading the dialogue, I realize that Claudia is trying to make amends with Tyler, but he's being defensive and hurtful. Claudia's tone is conciliatory, and she's trying to understand Tyler's perspective. Given this context, I should have scored Claudia's emotions differently. Next time, I will pay closer attention to the dialogue and consider the emotional dynamics between the characters. I will also consider the fact that Claudia is trying to apologize and make amends, which suggests that she is likely feeling regret and defensiveness, but not pride. I will also consider the fact that Tyler's hurtful comments are likely causing Claudia to feel regret and defensiveness, but not confusion.} \\
\\ 
\texttt{Trial \#2: New plan: I overestimated Claudia's emotional state and didn't consider the context of the conversation. Upon re{-}reading the dialogue, I realize that Claudia is trying to make amends with Tyler, but he's being defensive and hurtful. Claudia's tone is conciliatory, and she's trying to understand Tyler's perspective. Given this context, I should have scored Claudia's emotions differently. Next time, I will pay closer attention to the dialogue and consider the emotional dynamics between the characters. I will also consider the fact that Claudia is trying to apologize and make amends, which suggests that she is likely feeling regret and defensiveness, but not pride. I will also consider the fact that Tyler's hurtful comments are likely causing Claudia to feel regret and defensiveness, but not confusion.} \\
\\ 
\texttt{To improve my scoring, I will break down the dialogue into key moments and analyze the emotional cues in each moment. I will also consider the power dynamics between the characters and how they affect the emotional tone of the conversation. Specifically, I will look for moments where Claudia is trying to apologize or make amends, and moments where Tyler is being hurtful or defensive. By paying closer attention to these emotional dynamics, I can develop a more nuanced understanding of Claudia's emotional state and score her emotions more accurately.} \\
\\ 
\texttt{Regret: 9} \\
\texttt{Pride: 0} \\
\texttt{Confusion: 1} \\
\texttt{Defensiveness: 10} \\
\\ 
\texttt{Trial \#3: New plan: I overestimated Claudia's emotional state and didn't consider the context of the conversation. Upon re{-}reading the dialogue, I realize that Claudia is trying to make amends with Tyler, but he's being defensive and hurtful. Claudia's tone is conciliatory, and she's trying to understand Tyler's perspective. Given this context, I should have scored Claudia's emotions differently. Next time, I will pay closer attention to the dialogue and consider the emotional dynamics between the characters. I will also consider the fact that Claudia is trying to apologize and make amends, which suggests that she is likely feeling regret and defensiveness, but not pride. I will also consider the fact that Tyler's hurtful comments are likely causing Claudia to feel regret and defensiveness, but not confusion.} \\
\\ 
\texttt{To improve my scoring, I will break down the dialogue into key moments and analyze the emotional cues in each moment. I will also consider the power dynamics between the characters and how they affect the emotional tone of the conversation. Specifically, I will look for moments where Claudia is trying to apologize or make amends, and moments where Tyler is being hurtful or defensive. By paying closer attention to these emotional dynamics, I can develop a more nuanced understanding of Claudia's emotional state and score her emotions more accurately.} \\
\\ 
\texttt{I will also pay closer attention to the fact that Claudia's initial statement is positive and conciliatory ("Tyler, you look well. Teaching suits you"), which suggests that she is trying to make a good impression and start the conversation on a positive note. This is likely to be followed by a moment of tension or conflict, which will reveal more about Claudia's emotional state.} \\
\\ 
\texttt{Regret: 9} \\
\texttt{Pride: 0} \\
\texttt{Confusion: 1} \\
\texttt{Defensiveness: 10} \\
\\ 
\\ 
\\ 
\texttt{Your task is to predict the likely emotional responses of a character in this dialogue:} \\
\\ 
\texttt{Claudia: Tyler, you look well. Teaching suits you.} \\
\texttt{Tyler: Yeah, well, someone's gotta do the dirty work.} \\
\texttt{Claudia: I didn't mean it like that. I...} \\
\texttt{Tyler: I know what you meant. You always do, don't you?} \\
\texttt{Claudia: Not always, no. But I try to understand, at least.} \\
\texttt{Tyler: Understand? Like you understood when you left me alone?} \\
\texttt{Claudia: That's not fair, Ty. I had to...} \\
\texttt{Tyler: What? Chase your dreams? Leave your past behind?} \\
\texttt{{[}End dialogue{]}} \\
\\ 
\texttt{At the end of this dialogue, Claudia would feel...} \\
\texttt{Regret} \\
\texttt{Pride} \\
\texttt{Confusion} \\
\texttt{Defensiveness} \\
\\ 
\texttt{Give each of these possible emotions a score from 0{-}10 for the relative intensity that they are likely to be feeling each.} \\
\\ 
\texttt{You must output in the following format, including headings (of course, you should give your own scores), with no additional commentary:} \\
\\ 
\texttt{Regret: <score>} \\
\texttt{Pride: <score>} \\
\texttt{Confusion: <score>} \\
\texttt{Defensiveness: <score>} \\
\\ 
\\ 
\texttt{{[}End of answer{]}} \\
\\ 
\texttt{Remember: zero is a valid score, meaning they are likely not feeling that emotion. You must score at least one emotion > 0.} \\
\\ 
\texttt{Your answer:} \\
            }
        }
    }
     \\ \\

    \theutterance \stepcounter{utterance}  
    & \multicolumn{4}{p{0.6\linewidth}}{
        \cellcolor[rgb]{0.561,0.737,0.561}{
            \makecell[{{p{\linewidth}}}]{
                \texttt{\tiny{[P1$\rangle$GM]}}
                \texttt{Regret: 8} \\
\texttt{Pride: 0} \\
\texttt{Confusion: 2} \\
\texttt{Defensiveness: 9} \\
            }
        }
    }
    & & \\ \\

    \theutterance \stepcounter{utterance}  
    & & & \multicolumn{2}{p{0.3\linewidth}}{
        \cellcolor[rgb]{0.9,0.9,0.9}{
            \makecell[{{p{\linewidth}}}]{
                \texttt{\tiny{[GM$|$GM]}}
                \texttt{Regret: 8} \\
\texttt{Pride: 0} \\
\texttt{Confusion: 2} \\
\texttt{Defensiveness: 9} \\
            }
        }
    }
    & & \\ \\

    \theutterance \stepcounter{utterance}  
    & & & \multicolumn{2}{p{0.3\linewidth}}{
        \cellcolor[rgb]{0.9,0.9,0.9}{
            \makecell[{{p{\linewidth}}}]{
                \texttt{\tiny{[GM$|$GM]}}
                \texttt{Regret: 8} \\
\texttt{Pride: 0} \\
\texttt{Confusion: 0} \\
\texttt{Defensiveness: 6} \\
            }
        }
    }
    & & \\ \\

    \theutterance \stepcounter{utterance}  
    & & & \multicolumn{2}{p{0.3\linewidth}}{
        \cellcolor[rgb]{0.9,0.9,0.9}{
            \makecell[{{p{\linewidth}}}]{
                \texttt{\tiny{[GM$|$GM]}}
                \texttt{game\_result = LOSE} \\
            }
        }
    }
    & & \\ \\

    \theutterance \stepcounter{utterance}  
    & & \multicolumn{4}{p{0.6\linewidth}}{
        \cellcolor[rgb]{0.95,0.95,0.95}{
            \makecell[{{p{\linewidth}}}]{
                \texttt{\tiny{[P1$\leftrightarrow$P]}}
                \texttt{Post{-}episode Reflexion (not sent to GM):} \\
\texttt{New plan: I overestimated Claudia's emotional state and didn't consider the context of the conversation. Upon re{-}reading the dialogue, I realize that Claudia is trying to make amends with Tyler, but he's being defensive and hurtful. Claudia's tone is conciliatory, and she's trying to understand Tyler's perspective. Given this context, I should have scored Claudia's emotions differently. Next time, I will pay closer attention to the dialogue and consider the emotional dynamics between the characters. I will also consider the fact that Claudia is trying to apologize and make amends, which suggests that she is likely feeling regret and defensiveness, but not pride. I will also consider the fact that Tyler's hurtful comments are likely causing Claudia to feel regret and defensiveness, but not confusion.} \\
\\ 
\texttt{To improve my scoring, I will break down the dialogue into key moments and analyze the emotional cues in each moment. I will also consider the power dynamics between the characters and how they affect the emotional tone of the conversation. Specifically, I will look for moments where Claudia is trying to apologize or make amends, and moments where Tyler is being hurtful or defensive. By paying closer attention to these emotional dynamics, I can develop a more nuanced understanding of Claudia's emotional state and score her emotions more accurately.} \\
\\ 
\texttt{I will also pay closer attention to the fact that Claudia's initial statement is positive and conciliatory ("Tyler, you look well. Teaching suits you"), which suggests that she is trying to make a good impression and start the conversation on a positive note. This is likely to be followed by a moment of tension or conflict, which will reveal more about Claudia's emotional state.} \\
\\ 
\texttt{However, I will also consider the fact that Claudia's tone becomes less conciliatory and more defensive as the conversation progresses. This suggests that she is becoming increasingly uncomfortable and defensive in response to Tyler's hurtful comments. I will take this into account when scoring her emotions.} \\
\\ 
\texttt{Regret: 9} \\
\texttt{Pride: 0} \\
\texttt{Confusion: 1} \\
\texttt{Defensiveness: 10} \\
            }
        }
    }
    & \\ \\

\end{supertabular}
}

\end{document}
